% ============================================================
%  Chapter 2 Tutorial (Feynman-First)
%  Source: Chapter_2.pdf (Partha Pratim Das thesis)
%  Figures stored in: figures/
%  Figure naming: Parth_ch2_Fig2_[N]_page[XX].png
%  Workflow rules: see TUTORIAL_WORKFLOW.md
% ============================================================

\documentclass[12pt, a4paper]{article}

% ---- core packages ----
\usepackage[utf8]{inputenc}
\usepackage[T1]{fontenc}
\usepackage{amsmath, amssymb, amsthm}
\usepackage{bm}
\usepackage{siunitx}

% ---- figures ----
\usepackage{graphicx}
\graphicspath{{figures/}}
\usepackage{float}
\usepackage{caption}
\usepackage{subcaption}

% ---- layout ----
\usepackage[a4paper, top=2.5cm, bottom=2.5cm, left=3cm, right=2.8cm]{geometry}
\usepackage{parskip}
\setlength{\parskip}{6pt}
\setlength{\parindent}{0pt}

% ---- hyperlinks ----
\usepackage[colorlinks=true, linkcolor=blue!60!black, citecolor=blue!50!black]{hyperref}
\usepackage{bookmark}
\makeatletter
\pdfstringdefDisableCommands{%
    \def\vspace#1{}%
    \def\rule#1#2{}%
    \def\color#1{}%
    \let\leavevmode@ifvmode\relax%
    \def\,{ }%
}
\makeatother

% ---- styled boxes ----
\usepackage{tcolorbox}
\tcbuselibrary{skins, breakable}

% ---- tables ----
\usepackage{booktabs}
\usepackage{array}
\usepackage{longtable}

% ---- section formatting ----
\usepackage{titlesec}
\usepackage{xcolor}

\definecolor{parthblue}{RGB}{15, 58, 112}
\definecolor{parthorange}{RGB}{185, 72, 10}
\definecolor{parthgreen}{RGB}{10, 100, 42}
\definecolor{parthred}{RGB}{170, 25, 25}
\definecolor{lightblue}{RGB}{216, 232, 252}
\definecolor{lightorange}{RGB}{255, 244, 220}
\definecolor{lightgreen}{RGB}{216, 247, 228}
\definecolor{lightred}{RGB}{255, 228, 228}
\definecolor{lightgray}{RGB}{245, 245, 245}

\titleformat{\section}{\large\bfseries\color{parthblue}}{{\thesection.}}{0.5em}{}
            [\vspace{-4pt}\rule{\textwidth}{0.5pt}]
\titleformat{\subsection}{\normalsize\bfseries\color{parthblue!80!black}}{{\thesubsection}}{0.5em}{}

% ---- tcolorbox styles ----
\newtcolorbox{circuitbox}[1][Circuit Description]{
    colback=lightblue, colframe=parthblue, coltitle=white,
    fonttitle=\bfseries\small, title={#1},
    breakable, enhanced, left=4pt, right=4pt, top=4pt, bottom=4pt}

\newtcolorbox{statebox}[1][State Variables \& Parameters]{
    colback=lightorange, colframe=parthorange, coltitle=white,
    fonttitle=\bfseries\small, title={#1},
    breakable, enhanced, left=4pt, right=4pt, top=4pt, bottom=4pt}

\newtcolorbox{eqnbox}[1][Governing Equations]{
    colback=lightgreen, colframe=parthgreen, coltitle=white,
    fonttitle=\bfseries\small, title={#1},
    breakable, enhanced, left=4pt, right=4pt, top=4pt, bottom=4pt}

\newtcolorbox{physbox}[1][Physical Interpretation]{
    colback=lightgray, colframe=gray!60, coltitle=white,
    fonttitle=\bfseries\small, title={#1},
    breakable, enhanced, left=4pt, right=4pt, top=4pt, bottom=4pt}

\newtcolorbox{warningbox}[1][Watch Out]{
    colback=lightred, colframe=parthred, coltitle=white,
    fonttitle=\bfseries\small, title={#1},
    breakable, enhanced, left=4pt, right=4pt, top=4pt, bottom=4pt}

% ---- helpers ----
\newcommand{\figref}[1]{Figure~\ref{#1}}

% ---- figure entry macro ----
% Usage: \figentry{label}{filename-no-ext}{caption}{width}
\newcommand{\figentry}[4]{%
    \begin{figure}[H]
        \centering
        \includegraphics[width=#4\textwidth]{#2}
        \caption{#3}
        \label{#1}
    \end{figure}%
}

% ============================================================
\title{\textbf{Chapter 2 Tutorial: A Comprehensive Comparison of}\\[2pt]
       \textbf{Six-Phase and Three-Phase Drives}\\[6pt]
       \large Feynman-First, Derivation-First Notes (Grad Level)}
\author{}
\date{\today}
% ============================================================

\begin{document}
\maketitle
\tableofcontents
\clearpage

% ============================================================
\section{How to Use This Tutorial}
\label{sec:how_to_use}
% ============================================================

This document rewrites Chapter~2 of the thesis as a step-by-step tutorial.
Every major result is derived from elementary principles (phasors, KCL/KVL,
power balance, and PWM timing), while preserving the original notation.

\begin{warningbox}[Workflow Gate]
We write and revise \emph{one section or subsection at a time}.
You review and explicitly approve ("green signal"). Only then will the next part
be written.
\end{warningbox}

\begin{statebox}[How to Read (Feynman-First)]
For each result, we follow the same sequence:
\begin{enumerate}
    \item \textbf{Start from a figure:} what is visibly happening?
    \item \textbf{Name variables:} what does each symbol represent (and sign conventions)?
    \item \textbf{Derive:} use elementary tools (phasors, KCL/KVL, power balance, PWM timing).
    \item \textbf{Interpret:} what changes across three-phase, ASP, and SSP?
\end{enumerate}
At the end of every written section we add a \textbf{checkpoint question}:
"Explain in your own words what problem was solved and what the key takeaway is."
\end{statebox}

\clearpage

% ============================================================
\section{Lesson Plan (Chapter 2 Split)}
\label{sec:lesson_plan}
% ============================================================

We split Chapter~2 into lessons to keep each lesson self-contained and checkable.

\begin{enumerate}
    \item \textbf{Lesson 1: Configurations and Comparison Metrics}
    \item \textbf{Lesson 2: Torque Ripple}
    \item \textbf{Lesson 3: DC Bus Voltage Requirement}
    \item \textbf{Lesson 4: DC Bus Capacitor Requirement}
    \item \textbf{Lesson 5: CMV, Fault Tolerance, Efficiency, Conclusions}
\end{enumerate}

Each lesson will end with a short checkpoint: “explain in your own words.”

\clearpage

% ============================================================
\section{Lesson 1: Configurations and Comparison Metrics (Section 2.1)}
\label{sec:lesson1}
% ============================================================

\textit{This lesson defines what “three-phase”, “asymmetrical six-phase (ASP)”, and
“symmetrical six-phase (SSP)” PMSMs mean in space, and sets up the comparison
axes used throughout Chapter~2.}

\figentry{fig:2_1}{Parth_ch2_Fig2_1_page48}
{Figure 2.1: Spatial winding distribution of three-phase, asymmetrical, and symmetrical six-phase PMSMs.}{0.95}

\begin{circuitbox}[What Figure 2.1 Is Actually Saying]
Think of the stator as a circle. Each phase winding occupies a fixed direction in
space. When the rotor turns, the induced back-EMF in each phase depends on the
relative angle between the rotor flux and that winding direction.

Figure~2.1 shows only \textbf{where the windings sit in space}:
\begin{enumerate}
    \item \textbf{Three-phase PMSM (Fig.~2.1(a)).}
    Three windings (R, Y, B) are placed 120$^{\circ}$ apart (electrical).
    Their star point is labeled N1.

    \item \textbf{Six-phase PMSM = two three-phase sets (Fig.~2.1(b,c)).}
    One set is RYB and the other is UVW. Each set is internally 120$^{\circ}$ apart.
    The only design choice is the \emph{spatial shift} between the two sets.

    \item \textbf{ASP vs SSP.}
    ASP means the UVW set is shifted by 30$^{\circ}$ relative to RYB.
    SSP means it is shifted by 60$^{\circ}$.
\end{enumerate}
\end{circuitbox}

\begin{statebox}[Key Definitions (Notation Matches the Thesis)]
\textbf{RYB set:} a standard three-phase set with 120$^{\circ}$ phase displacement.

\textbf{UVW set:} a second three-phase set with 120$^{\circ}$ phase displacement.

\textbf{Asymmetrical Six-Phase (ASP) PMSM:} two three-phase sets (RYB and UVW)
separated by 30$^{\circ}$ in space (Fig.~2.1(b)).

\textbf{Symmetrical Six-Phase (SSP) PMSM:} two three-phase sets (RYB and UVW)
separated by 60$^{\circ}$ in space (Fig.~2.1(c)).

\textbf{Neutral points:} Fig.~2.1 labels the star points as N1 (for RYB) and N2
(for UVW). In practice, a six-phase machine may be used as:
\begin{enumerate}
    \item \textbf{2N connection:} N1 and N2 are kept separate (two neutrals).
    \item \textbf{1N connection:} N1 and N2 are tied together (single neutral).
\end{enumerate}
The 1N vs 2N choice changes which zero-sequence / circulating currents can flow,
which later changes DC-bus ripple and common-mode voltage behavior.
\end{statebox}

\begin{physbox}[Why the 30$^{\circ}$ vs 60$^{\circ}$ Shift Matters]
If the machine back-EMF and flux linkage were perfectly sinusoidal, many outcomes
would look similar across three-phase and six-phase.

But the considered PMSM has \textbf{odd-order harmonics} in back-EMF.
Those harmonics create torque ripple and can also inject low-frequency current
components. The spatial shift between RYB and UVW decides whether a given harmonic
from one set \textbf{adds} to the other set or \textbf{cancels}.

That single geometric fact is the root cause of several Chapter~2 results:
ASP cancels some harmonic torque ripple terms that do not cancel in SSP, while SSP
has stronger cancellation in some DC-bus ripple and common-mode voltage mechanisms
when carrier interleaving is used.
\end{physbox}

\begin{warningbox}[Two Common Confusions]
\textbf{1) “Six-phase” does not mean 60$^{\circ}$ between every adjacent phase.}
Each three-phase set is still 120$^{\circ}$ apart; ASP/SSP only sets the shift
between the two sets.

\textbf{2) Spatial angle is electrical.}
The 30$^{\circ}$/60$^{\circ}$ in Fig.~2.1 are electrical angles (after accounting for
pole pairs), because that is what aligns with back-EMF harmonics and torque.
\end{warningbox}

\subsection*{Comparison Axes Used in Chapter 2}

Chapter~2 compares three drive options (three-phase, ASPI+ASP-PMSM, SSPI+SSP-PMSM)
along these engineering axes:
\begin{enumerate}
    \item torque ripple,
    \item DC bus voltage requirement,
    \item DC bus capacitor requirement,
    \item common-mode voltage (CMV),
    \item fault tolerance (e.g., one-phase open circuit),
    \item inverter efficiency.
\end{enumerate}

\begin{statebox}[Checkpoint (Explain in Your Own Words)]
Looking at Fig.~2.1 only: what is the single difference between ASP and SSP?
Also: what does “1N vs 2N” physically mean in terms of which neutral points are
connected?
\end{statebox}

\clearpage

% ============================================================
\section{Lesson 2: Torque Ripple (Section 2.2)}
\label{sec:lesson2}
% ============================================================

\textit{This lesson compares torque ripple between three-phase, ASP, and SSP PMSMs. Torque ripple arises from back-EMF harmonics interacting with currents, causing unwanted vibrations and noise. We see how the 30° vs 60° spatial shift between RYB and UVW sets enables harmonic cancellation in ASP for certain orders.}

\figentry{fig:2_2}{Parth_ch2_Fig2_2_page49}{Figure 2.2: (a) Measured R phase back EMF of the SSP-PMSM. (b) The frequency spectrum of the back EMF at the speed of the motor (6750RPM or 450Hz). The back EMF contains all odd-order harmonics. Only one phase back EMF is shown. All other phases also have identical back EMF.}{0.95}

\begin{circuitbox}[Back-EMF Harmonics: What Figure 2.2 Shows]
PMSMs are not perfectly sinusoidal machines. The back-EMF (voltage induced when the rotor spins) has a fundamental (1st harmonic) plus higher odd harmonics (3rd, 5th, 7th, etc.).

Figure 2.2(a) shows the measured back-EMF waveform for the R phase of the SSP-PMSM at 6750 RPM (450 Hz fundamental). It's not a pure sine wave — there are distortions.

Figure 2.2(b) is the FFT (frequency spectrum): the fundamental at 450 Hz, plus 5th harmonic (2250 Hz), 7th (3150 Hz), etc. All odd multiples.

Why odd harmonics? In PMSMs with distributed windings and rotor saliency, even harmonics are absent because the flux linkage is symmetric. Only odd harmonics exist due to nonlinear saturation or winding imperfections.
\end{circuitbox}

\begin{physbox}[Torque Ripple: Why It Matters]
Torque ripple is the unwanted fluctuation in motor torque as the rotor spins. It causes vibrations, noise, and reduced efficiency. In electric vehicles, high ripple can lead to passenger discomfort or mechanical wear.

In PMSMs, torque ripple mainly comes from two sources:
\begin{enumerate}
    \item Cogging torque: due to stator/rotor tooth interaction (not considered here).
    \item Harmonic torque: interaction between back-EMF harmonics and stator currents.
\end{enumerate}

When currents are sinusoidal, but back-EMF has harmonics, the torque has ripple at the harmonic frequencies.

For example, 5th harmonic back-EMF × fundamental current → 5th harmonic torque ripple.
\end{physbox}

\begin{eqnbox}[Deriving Torque Ripple from Harmonics]
The electromagnetic torque in a PMSM is:

\[ T = \frac{3}{2} p (\psi_d i_q - \psi_q i_d) \]

But for harmonics, we expand the back-EMF as a Fourier series:

\[ e_a = \sum_{n=1,3,5,\dots} E_n \sin(n \theta) \]

where \(\theta\) is the electrical angle.

The torque ripple component from the nth harmonic is approximately:

\[ T_{ripple,n} \propto E_n I_{fundamental} \cos(n \delta) \]

where \(\delta\) is the spatial shift between the RYB and UVW sets.

For ASP (\(\delta = 30^\circ\)), certain harmonics (like 5th, 7th) have \(n \delta\) angles that lead to cancellation.

For SSP (\(\delta = 60^\circ\)), different cancellation patterns.
\end{eqnbox}

\figentry{fig:2_3}{Parth_ch2_Fig2_3_page51}{Figure 2.3: Torque ripple for the back-EMF of Fig.2.2(a) for (a) three-phase PMSM (b) ASP-PMSM (c) SSP-PMSM. The torque ripple of three-phase and SSP-PMSM are the same. The torque ripple due to the 5th and 7th harmonics flux linkage is zero in ASP-PMSM.}{0.95}

\begin{circuitbox}[Harmonic Cancellation in Multiphase Machines]
Figure 2.3 shows simulated torque ripple waveforms for the three configurations.

(a) Three-phase: baseline ripple.

(b) ASP: lower ripple because 5th and 7th harmonic torques cancel due to the 30° shift.

(c) SSP: similar to three-phase; the 60° shift doesn't cancel those harmonics as effectively.

In ASP, the RYB and UVW sets produce opposing torque ripples from 5th and 7th harmonics, summing to zero.
\end{circuitbox}

\begin{statebox}[Key Takeaway]
ASP PMSMs have lower torque ripple than three-phase or SSP because the 30° spatial shift causes 5th and 7th harmonic torque components to cancel between the two three-phase sets.

SSP and three-phase have identical ripple for these harmonics.
\end{statebox}

\begin{statebox}[Checkpoint (Explain in Your Own Words)]
What causes torque ripple in PMSMs? How does the spatial shift between RYB and UVW sets in six-phase machines affect torque ripple, particularly for the 5th and 7th harmonics?
\end{statebox}

\clearpage

% ============================================================
\section{Lesson 3: DC Bus Voltage Requirement (Section 2.3)}
\label{sec:lesson3}
% ============================================================

% TODO (write after approval):
% - Insert Fig. 2.4
% - Derive Vdc,min and Mmax for three-phase, ASPI, SSPI

\clearpage

% ============================================================
\section{Lesson 4: DC Bus Capacitor Requirement (Section 2.4)}
\label{sec:lesson4}
% ============================================================

% TODO (write after approval):
% - Insert Fig. 2.5–2.17 in the correct order
% - Derive 3rd-harmonic ripple model (KCL at split DC bus)
% - Derive switching-frequency ripple and Kcap,rpl interpretation

\clearpage

% ============================================================
\section{Lesson 5: CMV, Fault Tolerance, Efficiency, Conclusions (Sections 2.5–2.8)}
\label{sec:lesson5}
% ============================================================

% TODO (write after approval):
% - Insert Fig. 2.18–2.23 (CMV + interleaving)
% - Insert Table 2.2 and interpret (fault tolerance)
% - Insert Fig. 2.24–2.28 and Table 2.3–2.4 (efficiency)
% - Insert Fig. 2.29 (wrap-up tradeoff)

\clearpage

% ============================================================
\appendix
% ============================================================

% ============================================================
\section{Appendix 1: Why No Even Harmonics in PMSM Back-EMF?}
\label{app:even_harmonics}
% ============================================================

% TODO: Write proof and discussion, refer to textbooks if needed.

\clearpage

% ============================================================
\section{Appendix 2: Derivation of Electromagnetic Torque and Torque Ripple}
\label{app:torque_derivation}
% ============================================================

% TODO: Proof of torque equation, derivation of harmonic torque ripple.

\end{document}
